\chapter{Background}
\label{ch:background}

This chapter reviews the five strands of work the thesis builds on: in-context learning and what it learns (\autoref{sec:bg-icl}), the contest between priors and demonstrations (\autoref{sec:bg-priors}), distribution shift and shortcuts (\autoref{sec:bg-shift}), demonstration selection (\autoref{sec:bg-selection}), and the tools for looking inside a model (\autoref{sec:bg-mech}).
It closes with the gap this thesis addresses (\autoref{sec:bg-gap}).
\todo{Fold in the strongest material from the earlier literature review (Lit-review.pdf), and check each claim against the paper itself.}

\section{In-context learning}
\label{sec:bg-icl}

\citet{brown2020language} showed that a sufficiently large language model can perform a new task from a few demonstrations in its prompt, with no gradient updates; \citet{dong2024survey} survey the area since.
What the model learns from the demonstrations is contested.
\citet{min2022rethinking} found that replacing demonstration labels with random ones often costs surprisingly little, suggesting that demonstrations mainly convey the input distribution, the label space and the format.
\citet{pan2023what} separate \emph{task recognition}, retrieving a task the model already knows, from \emph{task learning}, acquiring an input--label mapping from the demonstrations themselves.
\citet{xie2022explanation} frame in-context learning as implicit Bayesian inference, in which demonstrations update a prior over latent concepts.
These views matter here because a model that mostly recognises tasks will follow its priors whenever the demonstrations disagree with them.

Controlled studies train or probe models on synthetic function classes.
\citet{garg2022what} show that transformers trained from scratch can learn linear functions and decision trees in context, and \citet{zhang2024trained} find that such in-context learners of linear functions are brittle under covariate shift.
Tabular foundation models are trained on synthetic task priors in the same spirit \citep{muller2022transformers,hollmann2023tabpfn,hollmann2025accurate}.
For pretrained language models, tabular rows are serialised as text \citep{hegselmann2023tabllm,dinh2022lift}, and classification is read from the probabilities of the label tokens \citep{brown2020language}, which are affected by surface-form competition \citep{holtzman2021surface} and by a bias toward particular labels that contextual calibration corrects \citep{zhao2021calibrate}.
Public tabular datasets can be memorised by large models \citep{bordt2024elephants}, one reason to prefer synthetic tasks.

\section{Priors against demonstrations}
\label{sec:bg-priors}

A model's prior is what it would predict with no demonstrations.
When feature names carry meaning, the prior is semantic: a feature called ``credit score'' suggests a direction before any example is seen.
\citet{dinh2022lift} and \citet{choi2022lmpriors} show that feature names change what language models do with tabular data, and \citet{nafar2025learning} compare named and anonymised features in in-context regression.

\citet{wei2023larger} tested what happens when the demonstrations contradict the prior, by flipping labels or replacing them with semantically unrelated ones: only large models follow the flipped labels, and smaller models stay with their priors.
At the scale studied here the picture is graded rather than all-or-nothing.

Llama-3.1-8B-Instruct agrees with fully flipped sentiment labels about half the time \citep{krishnakumar2025semantic}, and with flipped labels in-context accuracy plateaus below its level with the original labels \citep{kossen2024incontext}.

Related work on knowledge conflicts studies when a model trusts its context over its parametric memory \citep{longpre2021entity,xie2024adaptive}; the statistical notion of prior--data conflict goes back further \citep{evans2006checking}.
This thesis treats the question as a contest between prior and demonstrations whose outcome can be shifted, and asks whether the demonstrations themselves can be chosen to shift it.

\section{Distribution shift and shortcuts}
\label{sec:bg-shift}

\paragraph{Kinds of shift.}
Covariate shift changes the input distribution $P(X)$ while keeping $P(Y \mid X)$; label (prior-probability) shift changes $P(Y)$; concept shift changes $P(Y \mid X)$ itself \citep{shimodaira2000improving,quinonero2009dataset,morenotorres2012unifying}.
On real tabular benchmarks, most of the large performance drops are dominated by concept shift \citep{liu2023need}.
Benchmarks such as TableShift \citep{gardner2023tableshift} collect natural shifts, but each mixes several kinds, so the effect of one kind cannot be isolated.

\paragraph{Spurious features and shortcuts.}
A spurious feature is a feature with strong predictive power that predicts the label in the training environment for reasons that do not hold elsewhere.
Colored MNIST pairs a noisy causal signal with a colour whose correlation with the label reverses at test time \citep{arjovsky2019invariant}, and \citet{rosenfeld2021risks} model environmental features as class-conditional Gaussians.
Learners adopt such shortcuts readily \citep{geirhos2020shortcut,nagarajan2021understanding}.
Worst-group robustness is measured over (label, spurious attribute) groups \citep{sagawa2020distributionally}, and simple group balancing is a strong baseline \citep{idrissi2022simple}.
In-context learning inherits the problem: when demonstrations are tend to be consistent with several features, models show systematic preferences among them \citep{si2023measuring} and exploit shortcuts in the prompt \citep{tang2023large}, and in-context learners trained on tasks with spurious features adopt them \citep{harutyunyan2026incontext,chen2026incontext}.
Restricting a predictor to causal features is not automatically better: on tabular shifts, \citet{nastl2024causal} found predictors using all features generalised as well or better.

\section{Demonstration selection}
\label{sec:bg-selection}

\paragraph{Heuristic selection.}
The simplest heuristics retrieve the training examples nearest to the query \citep{liu2022what} or cover the input space with diverse examples \citep{levy2023diverse,gupta2023coverage}.
Others target robustness: comparable demonstrations, minimally edited to flip the label, break spurious associations \citep{fan2024comparable}, and misconfidence-based selection prefers examples on which the model's output departs from the true label \citep{xu2024misconfidence}.
Label balance and order are controlled because they bias predictions \citep{zhao2021calibrate,lu2022fantastically}.

\paragraph{Learned selection.}
Several methods train a selector from the language model's own feedback.
EPR \citep{rubin2022learning} and UDR \citep{li2023unified} train retrievers on the model's scores of candidate demonstrations, CEIL \citep{ye2023compositional} learns to choose whole sets, and \citet{zhang2022active} learn a selection policy by reinforcement learning.
Datamodels \citep{ilyas2022datamodels} regress a model's output on which training examples were present; \citet{chang2023data} fit them to in-context examples on the correct-class margin, and \citet{nguyen2023incontext} estimate influence from random subsets.
Signals from model internals and learned judges have also been proposed \citep{kato2025affinity,wang2026disp}; \citet{kato2025affinity} show that scores computed from induction-head representations track in-context accuracy but leave selection with them to future work.

\paragraph{Set-based architectures.}
The selector built in this thesis reads the whole pool of candidates at once.
It draws on In-Context Risk Minimization \citep{gupta2024context}, in which a transformer adapts to an environment by attending over other inputs from it, on the row tokens of TabPFN \citep{hollmann2023tabpfn}, and on attention over sets \citep{lee2019set}.

\section{Looking inside the model}
\label{sec:bg-mech}

The residual stream of a transformer \citep{vaswani2017attention} carries a running representation that every layer reads from and writes to \citep{elhage2021mathematical}.
The \emph{logit lens} reads an intermediate layer's prediction by applying the final normalisation and unembedding to it \citep{nostalgebraist2020logit}; its readings can be biased at early layers \citep{belrose2023eliciting}.
\citet{halawi2024overthinking} used it to find a \emph{critical layer} after which models given false demonstrations switch to the demonstrated labels.
Linear probes test whether information is linearly decodable from a layer \citep{alain2017understanding}, which shows that it is present, not that it is used \citep{belinkov2022probing}.
Steering adds a direction to the residual stream to change behaviour: activation addition \citep{turner2023activation} and contrastive activation addition \citep{rimsky2024steering} take differences of mean activations.
\citet{minder2025controllable} find a single direction that controls whether a model follows its context or its prior, and task and function vectors \citep{hendel2023incontext,todd2024function} compress an in-context task into one direction.

\section{The gap}
\label{sec:bg-gap}

Three gaps meet in this thesis.
First, selection methods are almost always evaluated in distribution, with no control over what kind of shift the queries undergo.
Second, the conflict between semantic priors and demonstrations has been studied by flipping \emph{labels}, not by changing what the \emph{features} mean, and not under shift.
Third, the learned selectors above are trained on the model's scores of whole prompts; none, as far as our search reached, uses the model's measured prior on each candidate, and none is tested on whether it helps a small model overcome that prior.
The nearest work, \citet{kato2025affinity}, is correlational and builds no selector.
This thesis builds the testbed that the first two gaps require, and uses it to test the third.
